\documentclass[10pt,a4paper]{amsart}
\usepackage[margin=19mm]{geometry}
\usepackage{amsmath,amssymb,mathtools}
\usepackage[scr=rsfs,cal=euler]{mathalfa}
\usepackage{xfrac}
\usepackage[unicode,hidelinks,bookmarks=false]{hyperref}
\newcommand{\ie}{i.e.\ }
\newcommand{\cf}{cf.\ }
\newcommand{\resp}{respectively\ }
\newcommand{\Subsection}{\subsection}
\providecommand{\nobreakdash}{\nobreak\hbox{-}}
\setlength{\emergencystretch}{2em}
\multlinegap0pt
\newtheorem{proposition}{Proposition}
\newcommand{\EE}{\mathbb{E}}
\title{On the Golomb--Dickman constant \\ under Ewens sampling}
\author{Jos\'e Ricardo G. Mendon\c{c}a, Luis Jehiel Negret}
\date{Source: arXiv:2603.23175v3 (CC BY 4.0)}
\begin{document}
\maketitle

\noindent\emph{Source and licence.} On the Golomb--Dickman constant \\ under Ewens sampling. Version arXiv:2603.23175v3 (CC BY 4.0), distributed by arXiv under the Creative Commons Attribution 4.0 International licence, http://creativecommons.org/licenses/by/4.0/, as recorded in the arXiv licence metadata for this exact version and shown on the abstract page https://arxiv.org/abs/2603.23175v3. Full licence notice: \texttt{licenses/arxiv-2603.23175-CC-BY-4.0.txt}.\par\smallskip

\noindent\textit{Editorial scope.} Counted unit: the article's proposition prop:smalltheta (source0.tex lines 285--292). The article's complete original proof of it is delivered with it (source0.tex lines 294--314, one \texttt{proof} environment of 21 lines). No mathematics is written, repaired or re-derived: the statement and the proof are the article's own lines, in the article's own notation, and no retained line is substituted or re-flowed.  Retained uncounted as the source-local context the proof needs: lines 255--258.  Nothing was omitted from any retained span, so the package carries no omission record.  No retained line contains a citation, so the wrapper carries no bibliography and no bibliography entry.  No retained line contains a figure environment, an image-inclusion command or drawing code, so no image is delivered and the package declares no asset dependency.  Editorial formatting only, in addition: an AMS \texttt{amsart} preamble that restates the article's own declarations and macros used by the retained lines, namely the environments \texttt{proposition} on separate counters, together with the standard packages those lines need.  The retained environments print in this document as Proposition 1 in order of appearance, whereas the article prints the same items with the numbering of its own full document.  The complete native source of the article as distributed in the arXiv e-print for this exact version is bundled unchanged as \texttt{sources/197059/source0.tex}.
\begin{equation}
\label{eq:ggd}
\lambda_{\theta} = \EE[Y_{\theta}] = \int_{0}^{\infty} \exp[-t-\theta E_{1}(t)]\,dt.
\end{equation}

\begin{proposition}
\label{prop:smalltheta}
As $\theta \to 0$,
\begin{equation}
\label{eq:smalltheta}
\lambda_{\theta} = 1 - (\ln{2})\theta + O(\theta^{2}).
\end{equation}
\end{proposition}

\begin{proof}
Expanding $\exp[-\theta E_{1}(t)] = 1-\theta E_{1}(t) + O(\theta^{2})$ in~\eqref{eq:ggd} and integrating term by term gives
\begin{equation}
\begin{split}
\lambda_{\theta} &= \int_{0}^{\infty} e^{-t}\,dt 
-\theta\int_{0}^{\infty} e^{-t} E_{1}(t)\,dt + O(\theta^{2}) \\
& = 1 - \theta \int_{0}^{\infty} e^{-t} E_{1}(t)\,dt + O(\theta^{2}).
\end{split}
\end{equation}
Substituting the definition of $E_{1}(t)$ and swapping the order of integration yields for the remaining integral above
\begin{equation}
\label{eq:firstorder}
\begin{split}
\int_{0}^{\infty} e^{-t} E_{1}(t)\,dt 
&= \int_{0}^{\infty} e^{-t} \int_{t}^{\infty} \frac{e^{-u}}{u}\,du\,dt \\
&= \int_{0}^{\infty} \frac{e^{-u}}{u} \int_{0}^{u} e^{-t}\,dt\,du 
= \int_{0}^{\infty} \frac{1 - e^{-u}}{u}\,e^{-u}\,du.
\end{split}
\end{equation}
If we write $(1-e^{-u})/u = \int_{0}^{1} e^{-uv}\,dv$ and integrate \eqref{eq:firstorder} in~$du$ first we get $\ln{2}$, and we are done.
\end{proof}

\end{document}
